\thispagestyle{empty}

\begin{center}
\vspace*{2.2cm}

{\sffamily\fontsize{34}{38}\selectfont\bfseries\color{ink} How a PGM Map\\[2pt] Gets Made}

\vspace{14pt}
{\color{rule}\rule{0.32\textwidth}{1.1pt}}
\vspace{14pt}

{\sffamily\large\color{ink60} The division of labour between an authoring agent\\ and \textsf{pgm-studio}}

\vspace{2.4cm}

\begin{minipage}{0.76\textwidth}
\small\color{ink}
Two maps are built from first principles in these pages: a Capture the Wool
board, composed by the studio and then adapted into a place, and a
destroy-the-monument board authored from nothing at all. Every stage of both is
shown as it came out of the machine. At each stage the paper asks the same
question. \emph{Who decided this?}
\end{minipage}

\vfill

{\footnotesize\sffamily\color{ink60}
Written from the running studio, September 2026\par}

\vspace{1.2cm}
\end{center}

\clearpage

% ---------------------------------------------------------------------------
\thispagestyle{empty}
\vspace*{1cm}

\section*{Purpose, Audience and Method}
\addcontentsline{toc}{chapter}{Purpose, Audience and Method}

\noindent
\textsf{pgm-studio} is a program for making Minecraft PvP maps of the kind a PGM
server loads. It is driven over HTTP, which means the thing operating it need not
be a person sitting at a screen. For most of the maps described here, the operator
was a language model with a terminal and a list of endpoints. This paper is an
account of what such a machine actually does, and of what the studio does for it.

\medskip\noindent
That question turns out to be the interesting one. A map is not produced by a
generator that takes a seed and returns a finished world, and it is not produced
by an author typing coordinates either. It is produced by two parties handing work
back and forth, and the line between them moves from stage to stage. The studio
arranges a board and the agent decides what the ground is. The agent states a
theme and the studio decides which block lands on which cell. The agent asks for
something, the studio refuses, and the refusal is where a good deal of the map
comes from. Naming that line precisely, stage by stage, is the whole of what
follows.

\medskip\noindent
The paper is in three parts, and they can be read in any order.

\begin{description}[leftmargin=1.6em,style=nextline,itemsep=5pt]
  \item[Part~I. The System] What a map is made of, the four documents that
    describe one, and the division of labour stated plainly. Three short
    chapters; read these first whatever else you read.
  \item[Part~II. The Instruments] One chapter per capability of the studio: the
    composer, the rules, shapes, relief, layers, terrain, houses, dressing, the
    export, reading a world back, and what an agent reads before it starts. This
    is the reference half, meant to be dipped into rather than read through.
  \item[Part~III. Two Maps, End to End] The two boards, every stage pictured, and
    a chapter on what an authoring session actually looks like from the outside,
    reconstructed from a recording of every call the agent made.
\end{description}

\medskip\noindent
A reader who plays these maps and wants to know where they come from can read
Part~I and then Part~III, and skip the middle entirely. A reader who wants to
drive the studio wants Part~II. Terms are defined where they are first used, and
collected again in the glossary at the back.

\medskip\noindent
One discipline is worth stating up front, because it constrains what this paper is
permitted to say. Nothing here claims that a map \emph{plays} a particular way
unless that claim traces to one of two sources: a ruling by the repository's
author, cited as such, or a measurement over the corpus of roughly 350 published
PGM maps, quoted with its number. The source code can say what the software does
and the corpus can say what mapmakers have done. Neither can say what is good.
Where the honest answer was that nobody had measured it, the claim is absent
rather than invented.

\clearpage
